\documentclass[10pt,a4paper]{amsart}
\usepackage[margin=19mm]{geometry}
\usepackage{amsmath,amssymb,mathtools}
\usepackage[scr=rsfs,cal=euler]{mathalfa}
\usepackage{xfrac}
\usepackage[unicode,hidelinks,bookmarks=false]{hyperref}
\newcommand{\ie}{i.e.\ }
\newcommand{\cf}{cf.\ }
\newcommand{\resp}{respectively\ }
\newcommand{\Subsection}{\subsection}
\providecommand{\nobreakdash}{\nobreak\hbox{-}}
\setlength{\emergencystretch}{2em}
\multlinegap0pt
\usepackage{amssymb}
\usepackage{mathtools}
\usepackage{float}
\usepackage[shortlabels]{enumitem}
\usepackage{tikz}
\newtheorem{lemma}{Lemma}
\title{The Exact End-Degree Threshold for\\ Finite-Width Directed Hexagonal Grids}
\author{Dan Zhu, Zhenhua Lyu}
\date{Source: arXiv:2608.21650v1 (CC BY 4.0)}
\begin{document}
\maketitle

\noindent\emph{Source and licence.} The Exact End-Degree Threshold for\\ Finite-Width Directed Hexagonal Grids. Version arXiv:2608.21650v1 (CC BY 4.0), distributed by arXiv under the Creative Commons Attribution 4.0 International licence, http://creativecommons.org/licenses/by/4.0/, as recorded in the arXiv licence metadata for this exact version and shown on the abstract page https://arxiv.org/abs/2608.21650v1. Full licence notice: \texttt{licenses/arxiv-2608.21650-CC-BY-4.0.txt}.\par\smallskip

\noindent\textit{Editorial scope.} Counted unit: the article's lemma lem:Gamma3 (source0.tex lines 472--474). The article's complete original proof of it is delivered with it (source0.tex lines 476--522, one \texttt{proof} environment of 47 lines). No mathematics is written, repaired or re-derived: the statement and the proof are the article's own lines, in the article's own notation, and no retained line is substituted or re-flowed.  No further source-local context is needed: every cross-reference of every retained line resolves inside the two delivered spans.  Nothing was omitted from any retained span, so the package carries no omission record.  No retained line contains a citation, so the wrapper carries no bibliography and no bibliography entry.  No retained line contains a figure environment, an image-inclusion command or drawing code, so no image is delivered and the package declares no asset dependency.  Editorial formatting only, in addition: an AMS \texttt{amsart} preamble that restates the article's own declarations and macros used by the retained lines, namely the environments \texttt{lemma} on separate counters, together with the standard packages those lines need.  The retained environments print in this document as Lemma 1 in order of appearance, whereas the article prints the same items with the numbering of its own full document.  The complete native source of the article as distributed in the arXiv e-print for this exact version is bundled unchanged as \texttt{sources/208265/source0.tex}.
\begin{lemma}\label{lem:Gamma3}
The digraph $\Gamma_3$ has an end of in-degree three, but contains no subdivision of $H_3$.
\end{lemma}

\begin{proof}
For $q\ge0$, write
\[
\begin{aligned}
v_{6q}&=r^1_{2q+1}, & v_{6q+1}&=r^2_{2q+1}, & v_{6q+2}&=r^2_{2q+2},\\
v_{6q+3}&=r^3_{2q+1}, & v_{6q+4}&=r^3_{2q+2}, & v_{6q+5}&=r^1_{2q+2},
\end{aligned}
\]
and set
\[
L_j=\{r^1_j,r^2_j,r^3_j\}.
\]
Each $L_j$ contains one vertex of each of $R_1,R_2,R_3$ and separates a finite initial segment of $\Gamma_3$
from all sufficiently late vertices. Hence every directed ray meets every sufficiently late $L_j$.

Let $Q_1,Q_2,Q_3$ be three pairwise disjoint rays. For every sufficiently large $j$, each $Q_i$ meets $L_j$, and the three intersections are distinct. Since $|L_j|=3$, $Q_1,Q_2,Q_3$ contains every vertex of every sufficiently late vertex of $L_j$, and therefore every sufficiently late vertex of $\Gamma_3$.

For sufficiently large $p$, let $x_p=1$ if the selected arc of $Q_1\cup Q_2\cup Q_3$ leaving $v_p$ is $v_pv_{p+5}$,
and let $x_p=0$ if it is $v_pv_{p+1}$. Set $x_p=0$ for odd $p$. Every sufficiently late vertex has selected in-degree one. At $v_q$ this gives
\[
(1-x_{q-1})+x_{q-5}=1,
\]
and hence
\begin{equation}
x_p=x_{p-4}
\end{equation}
for all sufficiently large $p$. Thus the tail pattern is determined by the two even residue classes modulo four. The corresponding successor map has the following numbers of cofinal orbits:
\[
\begin{array}{c|cccc}
(x_{0\pmod4},x_{2\pmod4})&(0,0)&(1,0)&(0,1)&(1,1)\\\hline
\text{number of cofinal orbits}&1&2&2&3.
\end{array}
\]
Indeed, tracing the successor map $p\mapsto p+1$ or $p\mapsto p+5$ in the four possible eventual residue patterns gives exactly the orbit counts displayed above. Since $Q_1,Q_2,Q_3$ are three cofinal orbits, the last case must occur. Hence $x_p=1$ for every sufficiently large even $p$. Up to permutation, the tails of $Q_1,Q_2,Q_3$ are the tails of $R_1,R_2,R_3$.

Suppose that $\Gamma_3$ contains a subdivision of $H_3$, and let $Q_1,Q_2,Q_3$ be its branch rays. By the preceding paragraph, their tails are the tails of $R_1,R_2,R_3$ and together contain every sufficiently late vertex of $\Gamma_3$. The vertex index strictly increases along every dipath in $\Gamma_3$. Choose a rung of a hypothetical $H_3$ subdivision whose endpoints lie beyond the finite initial part above. Every internal vertex of this rung is then also sufficiently late. Since every sufficiently late vertex already lies on one of the three branch rays and the interior of a rung is disjoint from all branch rays, every sufficiently late rung is a single arc.

The only late arcs between distinct rays have the cyclic orientations
\[
R_1\to R_2,
\qquad R_2\to R_3,
\qquad R_3\to R_1.
\]
Thus no pair of branch rays has sufficiently late rungs in both directions. This contradicts the definition of $H_3$. Therefore $\Gamma_3$ contains no subdivision of $H_3$.

Finally, for every ordered pair $R_i,R_j$, the periodic tail of $\Gamma_3$ contains directed $R_i$--$R_j$ linkages in infinitely many pairwise disjoint successive blocks, so $R_1,R_2,R_3$ lie in one end. Since every sufficiently late ray meets every sufficiently late three-vertex level $L_j$, and $|L_j|=3$, four pairwise disjoint rays are impossible. The end has in-degree exactly three.
\end{proof}

\end{document}
